\section{Validity Framework}

Validity is treated as a multidimensional evidence profile rather than a scalar score, following
the measurement principle that score interpretations require evidence for the intended use rather
than only high reliability or high accuracy \citep{messick1995validity,aera2014standards}. A
diagnostic can support a paper claim only when the artifact trail specifies the construct, the
response panel, the scoring rule, the diagnostic family, uncertainty or baseline checks where
available, and the claim boundary.

The framework separates three states that are often conflated. First, a benchmark can have a valid
response matrix for a diagnostic without proving the benchmark measures its intended construct.
Second, a diagnostic can flag suspicious items without proving those items are externally confirmed
errors. Third, a ranking can be sensitive to subjects or item subsets without becoming a true model
ranking. ValidEval records these distinctions in claims ledgers and evidence-state tables.

The active paper evidence uses this gate discipline: the 39-model MMLU panel passes the local
panel-validity gate, proxy IRT diagnostics are available, subject-level rank sensitivity is
artifact-backed, and MMLU-Redux direct/hash validation remains blocked.
